\documentclass[11pt,a4paper]{article}
\usepackage{auditstyle}
\lhead{Fresh Glimm--Jaffe proof audit}
\rhead{16 August 2026}
\newcommand{\eps}{\varepsilon}
\newcommand{\T}{\mathbb T}
\newcommand{\Z}{\mathbb Z}
\newcommand{\W}{\mathfrak W}
\newcommand{\PP}{\mathbb P}
\newcommand{\U}{\mathcal U}
\newcommand{\Tr}{\operatorname{Tr}}
\newcommand{\cth}{\operatorname{coth}}
\newcommand{\logN}{\mathcal L_N}

\title{Part II V4 in Lamport notation\\[3pt]
\large Fixed-time comparison, spectral transfer, ground-state limit, and assembly}
\author{Fresh reconstruction of \texttt{part\_ii/v4\_assembly.tex}}
\date{16 August 2026}

\begin{document}
\maketitle

\begin{resultbox}
This module proves the V4 assembly from the corrected V1--V3 modules.  In
particular, every use of V2's shifted estimate refers to the repaired
fundamental-band/alias proof, and every use of its fourth-chaos comparison includes
the four-edge term omitted by the source V2.  Thus the argument below validates the
claimed endpoint after those repairs; it does not retroactively validate the original
V2 proof.
\end{resultbox}
\tableofcontents

\section{Objects and corrected input contract}

Fix $L,m>0$, $\lambda\ge0$, a coarse scale $M$, and
$\xi=\eps_Mq+ip\in h_{M,L}$.  Let $N=M+r$, $\eps_N=2^{-r}\eps_M$, and
\[
 \logN=1+\log(1/\eps_N),\qquad
 \eta_D=2-\log_2 3>0.
\tag{0.1}
\]
The lattice Hamiltonian $K_{N,L}$ has ordered eigenvalues
$E_0(N)<E_1(N)\le\cdots$, simple ground vector $\Omega_N$, and ground state
$\omega_N$.  Put
\[
 e_0(N)=\frac12\sum_{k\in\Gamma_N}\gamma_N(k),\qquad
 \widetilde E_j(N)=E_j(N)-e_0(N),
\tag{0.2}
\]
\[
 D_N(t)=\sum_{j\ge1}e^{-t(E_j(N)-E_0(N))}.
\tag{0.3}
\]
The continuum operator $K_L$ constructed in the V3 module has eigenvalues
$\mu_0<\mu_1\le\cdots$, normalized positive ground vector $\Omega_L$, gap
$\gamma_L=\mu_1-\mu_0>0$, and
\[
 D_\infty(t)=\sum_{j\ge1}e^{-t(\mu_j-\mu_0)}<\infty.
\tag{0.4}
\]

Define the thermal Weyl functionals
\[
 F_N(t;\xi)=
 \frac{\Tr(e^{-tK_{N,L}}W_N(R_M^N\xi))}{\Tr e^{-tK_{N,L}}},\qquad
 F_\infty(t;\xi)=
 \frac{\Tr(e^{-tK_L}W_{\rm ct}(R_M^\infty\xi))}{\Tr e^{-tK_L}}.
\tag{0.5}
\]
Both have modulus at most one because the heat operator is positive trace class and
the Weyl operator is unitary.

The input contract is as follows.

\begin{enumerate}
\item V1 supplies the exact lattice marked-slice mean-shift formula, with the free
periodic Gaussian reference and with the interaction, not the full harmonic
potential, in the path weight.
\item The repaired V2 module supplies the common-noise laws, unshifted comparison,
the explicit four-edge correction, the exact fundamental-band/alias shifted
comparison, and uniform plain/shifted Nelson bounds.  Its standing time hypothesis is
$t_0L\ge1$.
\item V3 supplies the continuum trace, marked-Weyl formula, trace-class semigroup,
simple ground state, and positive gap.
\item The earlier free-rate and Weyl-reduction Lamport modules supply their stated
elementary estimates and $C^*$-algebraic extension theorem.
\end{enumerate}

No new printed theorem is imported in this module.

\section{Exact coupled representations}

Work on the repaired V2 common-noise probability space.  Let $\eta_N$ have the
lattice periodic law and $\eta$ the continuum periodic law.  Let
$\U_{N,t},\U_t$ be the corresponding interactions, and
$h^{(N)},h^{(\infty)}$ the deterministic jump interpolants.  Define
\[
 b_N(k)=\eps_M^{1/2}P_r(\eps_Mk)\widehat p_M(k),\qquad
 a_N(k)=\eps_M^{1/2}P_r(\eps_Mk)\widehat q_M(k),
\tag{1.1}
\]
and define $a_\infty,b_\infty$ by replacing $P_r$ with $P_\infty$.

\LS{1}{1}{Write both thermal functionals on this single coupling space.}

\LS{2}{1}{Define}
\[
 S_N=\frac1{8L}\sum_{k\in\Gamma_N}
 \gamma_N(k)\cth(t\gamma_N(k)/2)|b_N(k)|^2,
\quad
 S_\infty=\frac1{8L}\sum_{k\in\Gamma_\infty}
 \gamma(k)\cth(t\gamma(k)/2)|b_\infty(k)|^2,
\tag{1.2}
\]
\[
 \Phi_N(J_q)=\eps_N\sum_{x\in\Lambda_N}q_N(x)\eta_N(0,x),\qquad
 \Phi(J_q)=\int_{\T_L}q_\infty(x)\eta(0,x)dx.
\tag{1.3}
\]

\LS{2}{2}{Put}
\begin{align}
 \mathcal N_N&=\E[e^{i\Phi_N(J_q)}
 e^{-\U_{N,t}(\eta_N+h^{(N)})}],&
 \mathcal E_N&=\E[e^{-\U_{N,t}}],\tag{1.4}\\
 \mathcal N_\infty&=\E[e^{i\Phi(J_q)}
 e^{-\U_t(\eta+h^{(\infty)})}],&
 \mathcal E_\infty&=\E[e^{-\U_t}].\tag{1.5}
\end{align}

\LS{2}{3}{The exact representations are}
\[
 F_N=e^{-S_N}\frac{\mathcal N_N}{\mathcal E_N},
 \qquad
 F_\infty=e^{-S_\infty}\frac{\mathcal N_\infty}{\mathcal E_\infty}.
\tag{1.6}
\]

\Because The first is V1's mode-by-mode Mehler mean-shift identity with the
corrected interaction reference.  The second is equation (6.8) of the V3 module.
The repaired V2 covariance proposition identifies their Gaussian fields on one
space.  The Appell formula in V1/V3 shows that no extra Wick-subtraction term is
present.  All factors and normalizations appear explicitly in (1.1)--(1.5).

\LS{2}{4}{This proves $\langle1\rangle1$.}\QedStep

\section{Action and marked-source comparisons}

Throughout this section, $t\in[t_0,T]$, $t_0L\ge1$, and constants may depend on
the fixed coarse data but not on $N$ or $t$.  Put
$\bar w=\cth(t_0m/2)$.

\LS{1}{2}{Compare the action prefactors.}

\LS{2}{1}{$0\le S_N,S_\infty\le C_\xi$.}

\Because The D4 envelope gives
\[
 |b_\bullet(k)|^2\le C\eps_M\|\widehat p_M\|_\infty^2
 (1+\eps_M|k|)^{-2-\eta_D}.
\tag{2.1}
\]
Also $\gamma_N\le\gamma$, $\cth(t\gamma_N/2)\le\bar w$, and
$\sum_k(m+|k|)(1+\eps_M|k|)^{-2-\eta_D}<\infty$.

\LS{2}{2}{Let $\ell=\eps_M|k|$ and split at $\ell_r=2^{r/2}$.  On
$\ell\le\ell_r$,}
\[
 |\gamma_N\cth(t\gamma_N/2)-\gamma\cth(t\gamma/2)|
 \le C(t_0,m)\eps_N^2k^4,
\tag{2.2}
\]
\[
 ||b_N|^2-|b_\infty|^2|
 \le C_\xi2^{-r}\ell(1+\ell)^{-2-\eta_D}.
\tag{2.3}
\]

\Because For $w(y)=y\cth(ty/2)$,
$|w'(y)|\le\bar w+\sup_{x\ge t_0m/2}x/\sinh^2x$.  The lattice dispersion
estimate gives $0\le\gamma-\gamma_N\le\eps_N^2k^4/(24m)$, proving
(2.2).  The QMF telescope gives
$|P_r-P_\infty|\le C2^{-r}\ell(1+\ell)^{-1-\eta_D/2}$; multiply by
$|b_N|+|b_\infty|$ and use (2.1) to obtain (2.3).

\LS{2}{3}{Summing the low-mode dispersion term gives}
\[
 C4^{-r}\int_0^{\ell_r}\ell^4(1+\ell)^{-2-\eta_D}d\ell
 \le C2^{-r(1+\eta_D)/2}.
\tag{2.4}
\]

\Because The mode density in $\ell$ is constant.  At large $\ell$, the
integrand is $\ell^{2-\eta_D}$, whose integral is
$O(\ell_r^{3-\eta_D})$; multiply by $4^{-r}$.

\LS{2}{4}{Summing the low-mode filter term gives}
\[
 C2^{-r}\int_0^{\ell_r}\ell^2(1+\ell)^{-2-\eta_D}d\ell
 \le C2^{-r(1+\eta_D)/2}.
\tag{2.5}
\]

\Because The action weight grows at most linearly in $\ell$; (2.3) supplies the
other factor.  The integral is $O(\ell_r^{1-\eta_D})$.

\LS{2}{5}{The high-mode tails of both actions obey}
\[
 C\int_{\ell_r}^\infty\ell(1+\ell)^{-2-\eta_D}d\ell
 \le C2^{-r\eta_D/2}.
\tag{2.6}
\]

\LS{2}{6}{Therefore}
\[
 |S_N-S_\infty|\le C_\xi2^{-r\eta_D/2},\qquad
 |e^{-S_N}-e^{-S_\infty}|\le|S_N-S_\infty|.
\tag{2.7}
\]

\Because Add (2.4)--(2.6).  The last inequality follows from the mean-value
theorem because $S_N,S_\infty\ge0$ and $|(-e^{-x})|\le1$ there.

\LS{2}{7}{This proves $\langle1\rangle2$.}\QedStep

\LS{1}{3}{Compare the marked-slice source phases, including the asymmetric edge.}

\LS{2}{1}{For $r\ge1$, $a_N(\pm\pi/\eps_N)=0$.}

\Because The refinement mask contains the factor $m_0(\pm\pi)=0$ at the last
scale, so $P_r(\pm\pi2^r)=0$.  Thus the symmetrized V2 edge noise makes no
contribution to the linear marked source.  This fact applies only to this linear
source; it does not erase edge terms from covariance powers.

\LS{2}{2}{For $r\ge1$, orthogonality of the common Gaussian Fourier family gives}
\[
 \E[(\Phi_N-\Phi)^2]=\frac1{2Lt}\left[
 \sum_{\kappa\in D_N}
 |c_N(\kappa)^{1/2}a_N(-k)-c_\infty(\kappa)^{1/2}a_\infty(-k)|^2
 +\sum_{\kappa\in D_\infty\setminus D_N}c_\infty(\kappa)|a_\infty(-k)|^2
 \right].
\tag{2.8}
\]

\Because After Step $\langle2\rangle1$, the lattice source contains no edge
coefficient, so its Fourier noise is literally the restriction of the common family.
Subtract the two series and apply the Gaussian isometry.  The edge mode
$+\pi/\eps_N$ is included in the continuum-tail sum.

\LS{2}{3}{The filter-difference part of the first sum in (2.8) is at most}
\[
 C_\xi4^{-r}\int_0^{\pi2^r}\ell^2(1+\ell)^{-2-\eta_D}d\ell
 \le C_\xi2^{-r(1+\eta_D)}.
\tag{2.9}
\]

\Because
$t^{-1}\sum_{k_0}c_N(k_0,k)=\cth(t\gamma_N/2)/(2\gamma_N)\le\bar w/(2m)$;
square the QMF filter difference used in (2.3) and sum the spatial modes.

\LS{2}{4}{The covariance-symbol part is $O(\eps_N^2)$.}

\Because
$(\sqrt x-\sqrt y)^2\le|x-y|$ and the repaired V2 symbol estimate is
\[
 |c_N-c_\infty|(k_0,k)
 \le C\eps_N^2k^4(m^2+k_0^2+k^2)^{-2}.
\tag{2.10}
\]
The exact $k_0$ count gives
\[
 t^{-1}\sum_{k_0}|c_N-c_\infty|
 \le C\eps_N^2\frac{k^4}{(m^2+k^2)^{3/2}}
 \le C\eps_N^2(m+|k|).
\]
Multiplying by the D4 envelope for $|a_\infty|^2$ leaves a summable
$|k|^{-1-\eta_D}$ tail.

\LS{2}{5}{The missing-mode term in (2.8) is $O(2^{-r(1+\eta_D)})$.}

\Because Sum $t^{-1}\sum_{k_0}c_\infty\le\bar w/(2\gamma(k))$ against
$|a_\infty(k)|^2$ for $|k|\ge\pi/\eps_N$ and use the D4 tail.

\LS{2}{6}{After increasing the constant to cover the one fixed case $r=0$,}
\[
 \E[(\Phi_N-\Phi)^2]\le C_\xi2^{-r\eta_D/2},\qquad
 \|e^{i\Phi_N}-e^{i\Phi}\|_2\le C_\xi2^{-r\eta_D/4}.
\tag{2.11}
\]

\Because For $r\ge1$, combine (2.9)--(2.10) and the missing tail, then weaken
the exponent to a common one.  At $r=0$ the replacement of the edge noise in (2.8)
is not valid, but both variances are finite and bounded by a constant depending on
the fixed coarse scale, which equals the right side after enlarging $C_\xi$.  Finally
use $|e^{ix}-e^{iy}|\le|x-y|$.

\LS{2}{7}{This proves $\langle1\rangle3$.}\QedStep

\section{Numerators, denominators, and fixed-time convergence}

\LS{1}{4}{Bound the denominators away from zero.}

\LS{2}{1}{At the gate endpoint $(a,b)=(0,0)$,}
\[
 \E\U_{N,t}=3\lambda d_N(t)^2(2Lt),\qquad
 \E\U_t=3\lambda d_\infty(t)^2(2Lt).
\tag{3.1}
\]

\Because The ensemble-Wick chaoses of orders two and four are centered.  The Wick
recombination constant term is exactly $3d^2$; integrate the constant over volume
$2Lt$.

\LS{2}{2}{Since $0\le d_N,d_\infty\le\bar d$,}
\[
 \mathcal E_N,\mathcal E_\infty
 \ge e^{-C_0},\qquad C_0=3\lambda\bar d^2(2LT).
\tag{3.2}
\]

\Because Jensen's inequality for the convex function $e^{-x}$ gives
$\E e^{-\U}\ge e^{-\E\U}$; insert (3.1).

\LS{2}{3}{This proves $\langle1\rangle4$.}\QedStep

\LS{1}{5}{Compare the denominators.}

\LS{2}{1}{The repaired unshifted V2 comparison gives}
\[
 \|\U_{N,t}-\U_t\|_2
 \le C(1+t)(1+L)\eps_N\logN^{3/2}.
\tag{3.3}
\]

\Because Its fourth-chaos computation contains the one-edge and two-edge classes
and the nonzero four-edge class.  The last contributes the explicit positive term
\[
 \Delta_{4e}=4!(2Lt)^{-2}\frac32
 \sum_{k_{0,1}+\cdots+k_{0,4}=0}
 \prod_{i=1}^4 c_\times(k_{0,i},\pi/\eps_N),
\tag{3.4}
\]
bounded by
$9t\cth(2t_0)^3\eps_N^5/(256L^2)$.  Hence it lies below the common rate in
(3.3); it may not be deleted from the exact identity.

\LS{2}{2}{The plain lattice and continuum Nelson bounds and
$|e^{-x}-e^{-y}|\le|x-y|(e^{-x}+e^{-y})$ imply}
\[
 |\mathcal E_N-\mathcal E_\infty|
 \le C(1+t)(1+L)\eps_N\logN^{3/2}.
\tag{3.5}
\]

\Because Apply Cauchy--Schwarz to the product of the exponent difference and
$e^{-\U_N}+e^{-\U}$; (3.3) controls the first factor and the uniform $L^2$
Nelson bound the second.

\LS{2}{3}{This proves $\langle1\rangle5$.}\QedStep

\LS{1}{6}{Compare the numerators.}

\LS{2}{1}{Split}
\begin{align}
 |\mathcal N_N-\mathcal N_\infty|
 &\le\E[|e^{i\Phi_N}-e^{i\Phi}|
 e^{-\U_{N,t}(\eta_N+h^{(N)})}]\notag\\
 &\quad+\E[|e^{-\U_{N,t}(\eta_N+h^{(N)})}
 -e^{-\U_t(\eta+h^{(\infty)})}|].
\tag{3.6}
\end{align}

\LS{2}{2}{The first term in (3.6) is at most $C_\xi2^{-r\eta_D/4}$.}

\Because Apply Cauchy--Schwarz, (2.11), and the uniform shifted lattice
$L^2$ Nelson bound.

\LS{2}{3}{The repaired shifted V2 theorem gives}
\[
 \|\U_{N,t}(\eta_N+h^{(N)})-
 \U_t(\eta+h^{(\infty)})\|_2
 \le C_\xi\left(2^{-r(1+\eta_D)/4}
 +\eps_N\logN^{3/2}\right).
\tag{3.7}
\]

\Because Decompose it into the unshifted difference (3.3) and the difference of
the two shift increments.  The latter is proved in the repaired V2 module by:
(i) folding the lattice power symbols into the fundamental band; (ii) defining the
exact-output cross and continuum symbols; (iii) isolating the mixed alias term
$A_{\times,j}$; (iv) bounding filter, dispersion, high-output, and physical edge
remainders separately; and (v) applying the convolution telescope to every
deterministic coefficient power.  Thus no undefined $R_{\rm tail}$ is used in
(3.7).

\LS{2}{4}{The second term in (3.6) is bounded by the right side of (3.7).}

\Because Apply the exponential-difference inequality, Cauchy--Schwarz, (3.7),
and the paired uniform shifted Nelson bounds: V2 for finite $N$, V3 for the
continuum member.

\LS{2}{5}{Therefore}
\[
 |\mathcal N_N-\mathcal N_\infty|
 \le C_\xi(2^{-r\eta_D/4}+\eps_N\logN^{3/2}),
\tag{3.8}
\]
and
\[
 |\mathcal N_N|,|\mathcal N_\infty|\le B_1<\infty.
\tag{3.9}
\]

\Because The slower of the two powers of $2^{-r}$ is $\eta_D/4$.  Equation
(3.9) follows by deleting the unit-modulus phase and applying the $L^1$ shifted
Nelson bounds.

\LS{2}{6}{This proves $\langle1\rangle6$.}\QedStep

\LS{1}{7}{Prove fixed-time convergence of the interacting thermal functionals.}

\LS{2}{1}{The exact algebraic difference is}
\begin{align}
 F_N-F_\infty
 &= (e^{-S_N}-e^{-S_\infty})\frac{\mathcal N_N}{\mathcal E_N}
 +e^{-S_\infty}\frac{\mathcal N_N-\mathcal N_\infty}{\mathcal E_N}
 \notag\\
 &\quad+e^{-S_\infty}\mathcal N_\infty
 \frac{\mathcal E_\infty-\mathcal E_N}
 {\mathcal E_N\mathcal E_\infty}.
\tag{3.10}
\end{align}

\Because Put (1.6) over the common denominators and add/subtract the two
intermediate terms shown.

\LS{2}{2}{Equations (2.7), (3.2), (3.5), (3.8), and (3.9) give}
\[
 |F_N(t;\xi)-F_\infty(t;\xi)|
 \le C_\xi(2^{-r\eta_D/4}+\eps_N\logN^{3/2})\to0.
\tag{3.11}
\]

\LS{2}{3}{For every fixed $t\ge1/L$, (3.11) is admissible by taking
$t_0=\max(t/2,1/L)$ and $T=2t$.}

\Because Then $t_0\le t\le T$ and $t_0L\ge1$.  This is the genuine range of
the V2 comparison; no conclusion at arbitrarily small $t$ is inferred.

\LS{2}{4}{This proves $\langle1\rangle7$ (A6).}\QedStep

\section{Partition functions and exact Laplace transfer}

\LS{1}{8}{Compare the normal-ordered free partition functions.}

\LS{2}{1}{For every $t>0$,}
\[
 \widetilde Z_N^{(0)}(t):=Tr e^{-t(H_{N,L}^{(0)}-e_0(N))}
 =\prod_{k\in\Gamma_N}(1-e^{-t\gamma_N(k)})^{-1}.
\tag{4.1}
\]

\Because Subtracting $e_0(N)=\tfrac12\sum_k\gamma_N(k)$ removes the zero-point
energy from the product of one-oscillator traces.

\LS{2}{2}{For $k\in\Gamma_N$, the mean-value theorem gives}
\[
 |\log(1-e^{-t\gamma_N(k)})-
 \log(1-e^{-t\gamma(k)})|
 \le\frac{te^{-t\gamma_N(k)}}{1-e^{-tm}}
 |\gamma(k)-\gamma_N(k)|.
\tag{4.2}
\]

\Because The derivative of $y\mapsto\log(1-e^{-ty})$ is
$t/(e^{ty}-1)\le te^{-ty}/(1-e^{-tm})$ for $y\ge m$, and
$\gamma_N\le\gamma$.

\LS{2}{3}{Using}
\[
 0\le\gamma(k)-\gamma_N(k)\le\frac{\eps_N^2k^4}{24m},\qquad
 \gamma_N(k)\ge\max(m,2|k|/\pi),
\tag{4.3}
\]
the sum of (4.2) is $O(\eps_N^2)$.

\Because The remaining series
$\sum_k e^{-t\max(m,2|k|/\pi)}k^4$ converges with a bound independent of
$N$.

\LS{2}{4}{The omitted continuum modes contribute $O(e^{-2t/\eps_N})$.}

\Because The asymmetric complement begins at magnitude $\pi/\eps_N$, including
the positive edge.  Use
$|\log(1-e^{-x})|\le e^{-x}/(1-e^{-tm})$ and $\gamma(k)\ge|k|$; sum the
geometric spatial tail.

\LS{2}{5}{Hence}
\[
 |\log\widetilde Z_N^{(0)}(t)-\log Z_0(t)|
 \le C(m,L,t)\eps_N^2,\qquad
 \widetilde Z_N^{(0)}(t)\to Z_0(t).
\tag{4.4}
\]

\LS{2}{6}{This proves $\langle1\rangle8$.}\QedStep

\LS{1}{9}{Compare the interacting normal-ordered partition functions and obtain a
uniform lower ground-energy bound.}

Fix $t_*>0$ with $t_*L\ge1$.

\LS{2}{1}{For every $t\ge t_*$,}
\[
 \widetilde Z_N(t):=\Tr e^{-t(K_{N,L}-e_0(N))}
 =\widetilde Z_N^{(0)}(t)\mathcal E_N(t).
\tag{4.5}
\]

\Because Apply the finite-dimensional cyclic Mehler trace construction from V1/V3
to the lattice dispersions and interaction.  Normal ordering removes only the free
zero-point factor.  The repaired V2 proposition identifies the cyclic Gaussian law,
including the symmetrized edge column, and Wick recombination identifies the path
integrand with $\U_{N,t}$.

\LS{2}{2}{Equations (3.5), (4.4), and the V3 continuum partition identity imply}
\[
 \widetilde Z_N(t)\to
 Z_0(t)\mathcal E_\infty(t)=\Tr e^{-tK_L}=:Z_\infty(t)
 \quad(t\ge t_*).
\tag{4.6}
\]

\Because For this fixed $t$, use the admissible interval
$t_0=\max(t/2,1/L)$, $T=2t$ in (3.5).  This avoids the invalid choice
$t_0=t/2$ when $t<2/L$.

\LS{2}{3}{There is $c_\lambda<\infty$ such that}
\[
 \inf_N\widetilde E_0(N)\ge-c_\lambda.
\tag{4.7}
\]

\Because Fix any $t_1\ge t_*$.  Then
\[
 e^{-t_1\widetilde E_0(N)}\le\widetilde Z_N(t_1)
 =\widetilde Z_N^{(0)}(t_1)\mathcal E_N(t_1).
\]
The first factor is uniformly bounded by (4.4), and the second by the unshifted
$L^1$ Nelson bound.  Take logarithms and divide by $t_1$.

\LS{2}{4}{This proves $\langle1\rangle9$.}\QedStep

\LS{1}{10}{Prove the Laplace-transfer lemma in the exact range needed here.}

Let $\nu_N=\sum_{j\ge0}\delta_{x_j(N)}$ be locally finite eigenvalue point
measures, assume all $x_j(N)\ge c$, and suppose
\[
 \int e^{-tx}d\nu_N(x)\to\int e^{-tx}d\nu_\infty(x)
 \quad\hbox{for every }t\ge t_*.
\tag{4.8}
\]

\LS{2}{1}{The span $\mathscr A_{t_*}$ of $x\mapsto e^{-tx}$,
$t\ge t_*$, is uniformly dense in $C_0([c,\infty))$.}

\Because It is a self-adjoint algebra: products add the two parameters and remain
at least $t_*$.  It separates points and vanishes nowhere.  The locally compact
Stone--Weierstrass theorem applies.

\LS{2}{2}{For $f\in C_c([c,\infty))$, put $\widetilde f(x)=f(x)e^{t_*x}$ and
choose $g\in\mathscr A_{t_*}$ with
$\|\widetilde f-g\|_\infty<\delta$.  Then}
\[
 \left|\int f\,d\nu_N-\int ge^{-t_*x}\,d\nu_N\right|
 \le\delta\int e^{-t_*x}\,d\nu_N.
\tag{4.9}
\]

\Because Multiply the pointwise approximation error by the finite measure
$e^{-t_*x}d\nu_N$.  This weighted step is necessary: $\nu_N$ itself generally has
infinite total mass.

\LS{2}{3}{The right side of (4.9) is uniformly $O(\delta)$, and the middle
integrals converge.}

\Because Equation (4.8) at $t_*$ makes their masses uniformly bounded.  Since
$g$ is a finite linear combination of $e^{-tx}$ with $t\ge t_*$,
$ge^{-t_*x}$ is a finite linear combination with parameters at least $2t_*$, so
(4.8) applies.

\LS{2}{4}{Let $N\to\infty$ and then $\delta\downarrow0$ in (4.9).  Thus
$\nu_N\to\nu_\infty$ vaguely.}

\LS{2}{5}{Vague convergence implies convergence of every ordered atom with
multiplicity:}
\[
 x_j(N)\to x_j(\infty).
\tag{4.10}
\]

\Because Around the first $j+1$ limit atoms choose finitely many pairwise disjoint
intervals whose endpoints are not atoms and which contain no other limit atoms.
Continuous upper and lower approximations to their indicators show that the integer
masses in each interval eventually equal the limiting multiplicities.  Shrink the
intervals.  Induction in $j$ yields (4.10).

\LS{2}{6}{For $t\ge t_*$,}
\[
 D_N(t)=e^{tx_0(N)}\int e^{-tx}d\nu_N(x)-1
 \to D_\infty(t).
\tag{4.11}
\]

\Because Combine (4.8) and (4.10) for $j=0$.  No conclusion is asserted for
$t<t_*$; the hypotheses supply no control there.

\LS{2}{7}{This proves $\langle1\rangle10$.}\QedStep

\LS{1}{11}{Apply the transfer lemma to the interacting spectra.}

\LS{2}{1}{Take
$x_j(N)=\widetilde E_j(N)$ and $x_j(\infty)=\mu_j$.  Equations (4.6) and
(4.7) verify (4.8) and the common lower support bound.}

\LS{2}{2}{Therefore, for every fixed $j$,}
\[
 \widetilde E_j(N)\to\mu_j.
\tag{4.12}
\]

\LS{2}{3}{In particular,}
\[
 E_1(N)-E_0(N)	o\mu_1-\mu_0=\gamma_L>0,
\tag{4.13}
\]
and for every $t\ge t_*$,
\[
 D_N(t)\to D_\infty(t)<\infty.
\tag{4.14}
\]

\Because The zero-point shift cancels from each eigenvalue difference and from
$D_N$.  Apply (4.10)--(4.11) and the V3 strict gap.

\LS{2}{4}{This proves $\langle1\rangle11$ (A7).}\QedStep

\section{Thermal-to-ground-state interchange and the gate}

\LS{1}{12}{Prove the spectral tail estimate.}

Let $K$ have discrete spectrum $E_0<E_1\le\cdots$, normalized ground vector
$\Omega$, and let $W$ be unitary.

\LS{2}{1}{With $Z=e^{-tE_0}(1+D(t))$,}
\[
 \frac{\Tr(e^{-tK}W)}Z-\ip\Omega{W\Omega}
 =\frac1Z\sum_{j\ge1}e^{-tE_j}
 \left(\ip{\phi_j}{W\phi_j}-\ip\Omega{W\Omega}\right).
\tag{5.1}
\]

\Because Expand the trace in an orthonormal eigenbasis and cancel the $j=0$ term.

\LS{2}{2}{Since every displayed matrix element has modulus at most one,}
\[
 \left|\frac{\Tr(e^{-tK}W)}{\Tr e^{-tK}}-\ip\Omega{W\Omega}\right|
 \le\frac{2D(t)}{1+D(t)}\le2D(t).
\tag{5.2}
\]

\LS{2}{3}{This applies to every lattice $K_{N,L}$ and to $K_L$.}

\Because Lattice compactness and ground simplicity were proved in the model sheet;
V3 proves continuum trace class, discreteness, and ground simplicity.

\LS{2}{4}{This proves $\langle1\rangle12$.}\QedStep

\LS{1}{13}{Prove the abstract interchange of fixed-temperature and ground-state
limits.}

Suppose $a_N(t)$, $\ell_N$, and $D_N(t)$ obey
\[
 |a_N(t)-\ell_N|\le c_0D_N(t),\quad |a_N(t)|\le c_0,
\tag{5.3}
\]
$a_N(t)\to a_\infty(t)$ for each $t\ge t_0$,
$D_N(t_1)\to D_\infty(t_1)<\infty$ at one $t_1\ge t_0$, and
\[
 \liminf_N(E_1(N)-E_0(N))\ge\gamma>0.
\tag{5.4}
\]

\LS{2}{1}{For $t\ge t_1$,}
\[
 D_N(t)\le e^{-(t-t_1)(E_1(N)-E_0(N))}D_N(t_1).
\tag{5.5}
\]

\Because Factor $e^{-(t-t_1)(E_j-E_0)}$ from every summand and bound it by
the $j=1$ value.

\LS{2}{2}{There are $N_0,D_*<\infty$ such that, for all $N\ge N_0$,}
\[
 D_N(t)\le D_*e^{-(t-t_1)\gamma/2}.
\tag{5.6}
\]

\Because Convergence at $t_1$ bounds $D_N(t_1)$; the liminf in (5.4) gives an
eventual gap at least $\gamma/2$.

\LS{2}{3}{The continuum tail independently satisfies}
\[
 D_\infty(t)\le
 D_\infty(t_1)e^{-(t-t_1)\gamma_L}\to0.
\tag{5.7}
\]

\LS{2}{4}{Given $\delta>0$, choose $t$ so that the right sides of (5.6)--(5.7)
are below $\delta$.  Then}
\[
 |\ell_N-\ell_\infty|
 \le2c_0\delta+|a_N(t)-a_\infty(t)|.
\tag{5.8}
\]

\Because Add and subtract $a_N(t),a_\infty(t)$ and apply (5.3) to the two tails.

\LS{2}{5}{Let $N\to\infty$ in (5.8), then $\delta\downarrow0$.  Thus
$\ell_N\to\ell_\infty$.}

\LS{2}{6}{This proves $\langle1\rangle13$.}\QedStep

\LS{1}{14}{Close the Weyl gate at every fixed coarse scale.}

Set
\[
 t_*=\max(1,1/L),\quad
 a_N(t)=F_N(t;\xi),\quad
 \ell_N=\omega_N(W_N(R_M^N\xi)),
\tag{5.9}
\]
with the corresponding continuum definitions.

\LS{2}{1}{The bound (5.3) holds with $c_0=2$.}

\Because $|F_N|\le1$ and the spectral estimate (5.2) applies on both sides.

\LS{2}{2}{For each fixed $t\ge t_*$, $a_N(t)\to a_\infty(t)$.}

\Because $t_*\ge1/L$, so Step $\langle1\rangle7$ applies.

\LS{2}{3}{$D_N(t_*)\to D_\infty(t_*)<\infty$.}

\Because Step $\langle1\rangle11$ applies with the admissible lower time $t_*$.

\LS{2}{4}{$\liminf_N(E_1(N)-E_0(N))=\gamma_L>0$.}

\Because This is (4.13).

\LS{2}{5}{Step $\langle1\rangle13$ therefore gives}
\[
 \lim_{N\to\infty}
 \omega_N(\alpha_M^N(W_M(\xi)))
 =\omega^{\rm ct}_{L,\lambda}(\beta_M(W_M(\xi))).
\tag{5.10}
\]

\Because The embeddings are defined by
$\alpha_M^N(W_M(\xi))=W_N(R_M^N\xi)$ and
$\beta_M(W_M(\xi))=W_{\rm ct}(R_M^\infty\xi)$.

\LS{2}{6}{This proves $\langle1\rangle14$ (the gate).}\QedStep

\section{Inductive-limit state, identification, and symmetry}

\LS{1}{15}{Construct the continuum embedding of the inductive-limit algebra.}

Assume the already proved scaling identity
$\beta_N\alpha_M^N=\beta_M$ and injectivity of all maps.

\LS{2}{1}{On the algebraic union
$\bigcup_M\alpha_M^\infty(\W_M)$ define}
\[
 \beta(\alpha_M^\infty(A))=\beta_M(A).
\tag{6.1}
\]

\LS{2}{2}{Definition (6.1) is independent of the representative.}

\Because If $M\le M'$ and
$\alpha_M^\infty(A)=\alpha_{M'}^\infty(A')$, injectivity of
$\alpha_{M'}^\infty$ gives $A'=\alpha_M^{M'}(A)$.  Hence
$\beta_{M'}(A')=\beta_{M'}\alpha_M^{M'}(A)=\beta_M(A)$.

\LS{2}{3}{$\beta$ is isometric on the union.}

\Because An injective $*$-homomorphism between $C^*$-algebras is isometric;
thus
$\|\beta(\alpha_M^\infty(A))\|=\|\beta_M(A)\|=\|A\|
=\|\alpha_M^\infty(A)\|$.

\LS{2}{4}{It extends uniquely to an injective $*$-homomorphism}
\[
 \beta:\W_\infty\longrightarrow\W_{\rm ct},\qquad
 \beta\alpha_M^\infty=\beta_M.
\tag{6.2}
\]

\Because The algebraic union is dense in its inductive-limit completion.  An
isometry extends uniquely; multiplication and involution pass by continuity, and an
isometry is injective.

\LS{2}{5}{This proves $\langle1\rangle15$.}\QedStep

\LS{1}{16}{Apply the previously proved Weyl extension criterion.}

Define
\[
 \chi_M(\xi)=\lim_{N\to\infty}
 \omega_N(\alpha_M^N(W_M(\xi))).
\tag{6.3}
\]

\LS{2}{1}{The limit exists and equals}
\[
 \chi_M(\xi)=
 \omega^{\rm ct}_{L,\lambda}(\beta_M(W_M(\xi))).
\tag{6.4}
\]

\Because This is (5.10).

\LS{2}{2}{For each $M$, there is a unique state $\rho_M$ on $\W_M$ satisfying
$\rho_M(W_M(\xi))=\chi_M(\xi)$, and}
\[
 \omega_N(\alpha_M^N(A))\to\rho_M(A)\qquad(A\in\W_M).
\tag{6.5}
\]

\Because The Weyl-reduction module proves: pointwise convergence on all Weyl
generators, uniform state norm one, and density of finite Weyl linear combinations
extend the limit uniquely to the whole $C^*$-algebra.  Equation (6.4) also supplies
positive definiteness and continuity directly as a restriction of a state.

\LS{2}{3}{The family is consistent:}
\[
 \rho_N\circ\alpha_M^N=\rho_M.
\tag{6.6}
\]

\Because Check it on $W_M(\xi)$ using
$\alpha_M^N(W_M(\xi))=W_N(R_M^N\xi)$ and the compatibility of the continuum
refinement; extend by norm density.

\LS{2}{4}{There is a unique state $\Omega_\infty$ on $\W_\infty$ with}
\[
 \Omega_\infty\circ\alpha_M^\infty=\rho_M,\qquad
 \Omega_\infty=\omega^{\rm ct}_{L,\lambda}\circ\beta.
\tag{6.7}
\]

\Because Consistency defines a norm-one positive functional on the dense algebraic
union, hence a unique state on the completion.  Equations (6.2), (6.4) show that it
agrees there with the continuum pullback; density gives the second identity.

\LS{2}{5}{All states in (6.5)--(6.7) are invariant under the field reflection
$\vartheta(W(\zeta))=W(-\zeta)$.}

\Because Every refinement map is linear, hence intertwines $\vartheta$.  The
lattice Hamiltonian is even and commutes with its reflection unitary.  Ground-state
simplicity forces that unitary to multiply $\Omega_N$ by a phase, so
$\omega_N\circ\vartheta_N=\omega_N$.  Pass to the limits in (6.5), then use
consistency.  Equivalently, the earlier Weyl-reduction symmetry proposition gives
the same conclusion.

\LS{2}{6}{This proves $\langle1\rangle16$ and the full fixed-coarse theorem.}
\QedStep

\section{V4 conclusion and validity status}

\begin{resultbox}
Step $\langle1\rangle7$ proves fixed-$t$ thermal convergence with rate.
Steps $\langle1\rangle8$--$\langle1\rangle11$ prove normal-ordered partition
convergence, eigenvalue convergence, the eventual lattice gap, and tail convergence.
Steps $\langle1\rangle12$--$\langle1\rangle14$ interchange the thermal and ground
limits and prove the gate.  Steps $\langle1\rangle15$--$\langle1\rangle16$
identify the unique projective limit with the continuum vacuum and carry the
$\Z_2$ symmetry.

This conclusion is conditional on replacing the original V2 Lemma 6.4/Proposition
6.5 proof by the exact repaired comparison in the preceding Lamport module.  The
original Part II files, without that replacement and the four-edge correction, do
not prove their final theorem as written.
\end{resultbox}

\end{document}
